\section{Introduction}
\label{sec:introduction}

When a population contains only a few individuals, its extinction, persistence and
time spent near zero are distributional questions rather than small corrections to a
deterministic mean.  The models used later in this thesis therefore describe integer
counts that change at random event times.  Birth, death, transfer, absorption and
rupture are treated as abstract jump mechanisms, while the population count records
the state on which those mechanisms act.

This chapter develops two connected strands of the required probability.  The first
begins in discrete time with a binary Galton--Watson process.  Its simplicity makes
the distinction between mean behaviour and survival particularly sharp.  Below the
critical reproduction probability the population becomes extinct almost surely, but
the population conditional on survival approaches a non-degenerate limiting law.
For the binary model, the limiting conditional mean is the reciprocal of the
constant
\begin{equation}
  A(p)=\lim_{n\to\infty}\frac{S_n}{(2p)^n},
  \qquad 0<p<\tfrac12,
  \label{eq:intro-A}
\end{equation}
where \(S_n\) is the probability of survival to generation \(n\).  The value at the
left endpoint will always mean the right limit, \(A(0):=\lim_{p\downarrow0}A(p)\).
The constant in \cref{eq:intro-A} measures the cumulative effect of the nonlinear
part of the survival recursion before its late geometric decay becomes dominant.
It has an exact product representation, an exact series, useful two-sided bounds and
a controlled near-critical asymptotic.  It also has a less obvious interpretation as
a value of a Koenigs linearising function.  These facts explain both how \(A(p)\) can
be computed and why elementary curve fitting has proved unreliable.

The second strand concerns continuous-time jump processes and their generating
functions.  Once a process has more than one count, its forward equations form a
large or infinite coupled system.  A multivariate probability generating function
compresses that system into a first-order partial differential equation.  The method
of characteristics then converts the PDE into ordinary differential equations along
curves in the enlarged state--time space.  We develop this method through three
transfer models.  The first conserves particles, the second permits death after
transfer, and the third adds linear birth.  The final model has an infinite state
space and leads to a closed generating function involving a Gauss hypergeometric
function.  The calculation supplies the analytical template used later for
birth--death--catastrophe processes.

The two strands meet in several places.  Both use generating functions to replace a
countable family of probabilities by one analytic object.  Both are organised by
absorption: extinction in the branching process and transfer, death or rupture in
the continuous-time models.  Most importantly, both distinguish an unconditional
population, which may decay to zero, from a population observed under a survival or
non-absorption condition.  Quasi-stationarity is the language for that distinction.

The order of the chapter follows the mathematical dependence rather than the later
applications.  \Cref{sec:ctmc-pgf} gives only the continuous-time and generating-
function background needed for the calculations that follow.  The Galton--Watson
process, its extinction threshold and its critical lifetime are developed in
\cref{sec:galton-watson}.  \Cref{sec:qsd} introduces quasi-stationarity and obtains
the limiting conditional moments.  Founding populations, a discrete rupture
mechanism, the linear continuous-time birth--death process and the killed-generator
form of a catastrophe mechanism are treated in \cref{sec:small-populations}.  The
analysis of \(A(p)\) occupies
\cref{sec:constant-A}; the main text contains the exact results and proofs, while the
exploratory closed-form search is retained in a chapter-local appendix.  Finally,
\cref{sec:characteristics} derives and solves the generating-function PDEs.  Two
further appendices collect a hypergeometric integral identity and the extraction of
state probabilities from multivariate generating functions.

Standard background is stated concisely.  Proofs are given when the result is
specific to the models developed here, and numerical evidence is labelled as such.
This separation matters near criticality, where a visually accurate approximation
to \(A(p)\) can still produce a poor approximation to the conditional mean \(1/A(p)\).
